\documentclass{article}

\usepackage[preprint]{neurips_2026}
\usepackage[numbers]{natbib}
\usepackage[utf8]{inputenc}
\usepackage[T1]{fontenc}
\usepackage{hyperref}
\usepackage{url}
\usepackage{booktabs}
\usepackage{amsfonts}
\usepackage{amsmath}
\usepackage{microtype}
\usepackage{graphicx}
\usepackage{xcolor}
\usepackage{multirow}
\usepackage{siunitx}

\title{BrazilBench: A Reproducible Multi-Regime Benchmark for\\
Trading Agents on Brazilian Equities}

\author{
  Rafael D. Coelho \\
  Independent Researcher \\
  Fortaleza, CE, Brazil \\
  \texttt{rafael.d.cooelho@gmail.com} \\
}

\begin{document}

\maketitle

% -------------------------------------------------------
\begin{abstract}
We present BrazilBench, a reproducible benchmark for evaluating
autonomous trading agents on Brazilian equity markets
(\url{https://github.com/RAFAELDCOELHO/autonomous-hedge-fund}).
Existing benchmarks focus almost exclusively on U.S.\ markets, ignoring
structural differences of emerging economies: high interest rates,
currency volatility, and sharper crashes. BrazilBench defines four
macroeconomic regimes on B3 (bull\_2019, crisis\_2020,
recovery\_2021, high\_rates\_2022), a uniform agent interface that
carries macro context (SELIC, IPCA, USD/BRL) alongside prices, and a
simulation engine with SELIC as the country-appropriate risk-free
rate. This draft reports classical results only: five rule-based
agents (Buy-and-Hold, MACD, SMA crossover, Momentum, Random) on five
core B3 assets, a RandomAgent null distribution over $N{=}100$ seeds,
and equal-weight portfolio baselines. The null distribution and
portfolio tables are regenerated byte-for-byte by \texttt{make
reproduce} from committed price fixtures, with no API key and no
network access. Findings: \emph{(i)} regime sensitivity dominates
strategy differences --- MACD avoided the COVID-19 crash via a single
well-timed exit/re-entry on PETR4 ($+26.2\%$ vs.\ $-28.0\%$ B\&H;
3 fills), while Momentum lost up to $-28.5\%$ in the high-rates
regime from transaction-cost compounding; \emph{(ii)} an uninformed
random policy has a wide noise floor (per-cell standard deviation of
10--25\,pp), so single-seed baselines are noisy point estimates;
\emph{(iii)} the choice of risk-free rate (SELIC vs.\ Fed Funds)
changes Sharpe ratios by a factor of 2--3 and must be
country-specific. No LLM agent is evaluated in this draft: LLM
evaluation requires paid or GPU inference, is outside the
zero-cost reproduce path, and is left as future work.
\end{abstract}

% -------------------------------------------------------
\section{Introduction}
% -------------------------------------------------------

Recent years have seen a surge in LLM-powered trading frameworks, yet
their evaluation has remained almost exclusively focused on U.S.\ markets.
No standardized benchmark existed for testing trading agents on Brazilian
equities, motivating the creation of BrazilBench.

Brazil presents structural challenges that U.S.-centric benchmarks cannot
capture. The SELIC rate --- Brazil's benchmark interest rate --- fluctuated
between 2.0\% and 13.75\% during our study period, fundamentally altering
the opportunity cost of capital. The Brazilian real depreciated
significantly against the dollar over the same period. Most strikingly,
the Ibovespa index fell 45\% peak-to-trough during the COVID-19 crash
of 2020, nearly double the S\&P 500's decline. These are not exceptional
events --- they represent the normal operating environment of an emerging
market.

We introduce BrazilBench, a reproducible benchmark for evaluating
autonomous trading agents on Brazilian equity markets. Our contributions
are: (1) a standardized evaluation framework for trading agents on
B3-listed assets, with a macro-augmented agent state (SELIC, IPCA,
USD/BRL) available in the interface, four macroeconomic regimes
covering $2019$--$2022$, and SELIC as the risk-free rate; (2) a
zero-cost, CI-pinned reproduction path: \texttt{make reproduce}
regenerates the RandomAgent null distribution ($N{=}100$ seeds per
cell) and the equal-weight portfolio baselines from committed price
fixtures, with no API key or network access; and (3) regime-level
findings on classical strategies --- a single-event MACD case study
in the COVID-19 crash, Momentum whipsaw under high rates, the
dominance of Buy-and-Hold in sustained trends, and the effect of the
risk-free rate convention on cross-market Sharpe comparisons.

This draft deliberately reports \textbf{no LLM-agent results}. The
agent interface is designed so that LLM agents (cloud or open-weight)
can be evaluated on the same regimes, but such runs require paid API
calls or local GPU inference, are not part of \texttt{make reproduce},
and are not claimed here (Section~\ref{sec:limitations}).

The remainder of this paper is organized as follows. Section~2 reviews
related work. Section~3 describes our methodology. Section~4 presents
experimental results. Section~5 discusses key findings. Section~6
acknowledges limitations and outlines future work.

% -------------------------------------------------------
\section{Related Work}
% -------------------------------------------------------

\textbf{LLM trading frameworks.}
TradingAgents \cite{tradingagents2024} introduced a multi-agent LLM
framework where specialized analysts --- fundamentals, technicals,
sentiment --- collaborate through a structured graph to produce trading
decisions, demonstrating strong results on U.S.\ equities in 2023--2024.
FinAgent \cite{finagent2024} extended this paradigm with a
memory-augmented LLM that retrieves historical analogues before each
decision, improving robustness on unseen market conditions. BrazilBench
adopts the same agent interface philosophy but shifts the focus from
agent architecture to evaluation methodology, applying it to Brazilian
markets for the first time.

\textbf{Financial time-series foundation models.}
Kronos \cite{kronos2025} (Shi et al., 2026) is a decoder-only foundation model pre-trained
on over 12 billion K-line records from 45 global exchanges using a
domain-specific OHLCV tokenizer, achieving a 93\% improvement in RankIC
over leading time-series foundation models.
Chronos \cite{chronos2024} (Ansari et al., Amazon, 2024) is a
general-purpose time-series foundation model pre-trained on a large
corpus of diverse time series using language model objectives. Both
Kronos and Chronos target \textit{forecasting accuracy}. BrazilBench
targets a different axis: \textit{does a trading agent generate
risk-adjusted excess return over the SELIC rate across structurally
distinct market regimes?} A model with high forecasting accuracy can
still lose money once transaction costs, position sizing, and regime
shifts are introduced --- the gap between forecast accuracy and
tradeable alpha is precisely what BrazilBench is designed to measure.
Evaluating Kronos and Chronos as comparator agents is left as future
work.

\textbf{LLM financial language tasks.}
FinGPT \cite{fingpt2023} addresses LLM fine-tuning for financial
language tasks such as sentiment analysis and earnings prediction.
These works evaluate language model quality on text, not trading agent
decisions. BrazilBench is orthogonal: we evaluate what agents
\textit{do} in markets, not what they \textit{say} about them.

\textbf{Trading simulation and backtesting frameworks.}
Backtrader \cite{backtrader2015} and QuantConnect's Lean engine
\cite{quantconnect2013} are widely used open-source backtesting
platforms supporting multi-asset strategies, slippage models, and live
trading integration. BrazilBench differs from these platforms in two
ways: it defines standardized market regimes as evaluation units rather
than arbitrary date ranges, and it natively integrates macroeconomic
indicators from the Brazilian Central Bank \cite{bacen} as agent inputs.

\textbf{Market regime detection.}
Hamilton \cite{hamilton1989} introduced the hidden Markov model for
detecting shifts between economic regimes, a methodology widely applied
to equity markets. Ang and Bekaert \cite{angbekaert2002} characterized
international stock return regimes and showed that regime-conditional
Sharpe ratios differ substantially across countries --- a finding
directly motivating our multi-regime design. The four regimes in
BrazilBench are defined by macroeconomic events rather than statistical
detection, making them interpretable and replicable.

\textbf{Transaction costs and emerging markets.}
Lesmond et al.\ \cite{lesmond1999} documented that transaction costs in
emerging markets are substantially higher than in developed markets and
can eliminate apparent alpha. Adler and Dumas \cite{adlerdumas1983}
showed that currency risk introduces a separate risk factor for
international investors not captured by U.S.-only models --- directly
motivating our inclusion of USD/BRL as a macro feature. To our
knowledge, no prior benchmark evaluates autonomous trading agents on
B3-listed assets with SELIC as a country-appropriate risk-free rate.

% -------------------------------------------------------
\section{Methodology}
% -------------------------------------------------------

\subsection{Market Regimes}

We selected four periods that represent structurally distinct episodes
in Brazilian market history: the 2019 bull run driven by pension reform
optimism, the COVID-19 crash of 2020, the 2021 commodity-led recovery,
and the 2022 high-rates period under SELIC at 13.75\%. We label the
2022 period \texttt{high\_rates\_2022} rather than ``bear market''
because the Ibovespa closed $+5.6\%$ nominally. The bear
characterization applies only in risk-adjusted terms: with SELIC at
11.75\% as the opportunity cost of capital, the regime delivered a
Sharpe of $-0.19$, meaning a passive equity investor destroyed
risk-adjusted value relative to government bonds. ``Low
risk-adjusted return regime'' would be equally precise; we retain the
shorthand \texttt{high\_rates\_2022} throughout. Each regime captures
a different macroeconomic context, maximizing the diversity of
conditions under which agents are evaluated.

\begin{table}[h]
\centering
\caption{Brazilian market regimes used in BrazilBench.}
\label{tab:regimes}
\begin{tabular}{llllrp{4cm}}
\toprule
\textbf{Regime} & \textbf{Start} & \textbf{End} & \textbf{Days}
  & \textbf{Return} & \textbf{Context} \\
\midrule
bull\_2019        & 2019-01-02 & 2019-12-31 & 248 & +27.1\%
  & Post-pension reform, SELIC cuts \\
crisis\_2020      & 2020-02-03 & 2020-05-29 &  80 & $-23.9\%$
  & COVID-19 crash and recovery start \\
recovery\_2021    & 2021-01-04 & 2021-06-30 & 122 & +6.7\%
  & Vaccine rally, commodity surge \\
high\_rates\_2022 & 2022-01-03 & 2022-12-30 & 250 & +5.6\%
  & SELIC 13.75\%, Sharpe $-0.19$, election uncertainty \\
\bottomrule
\end{tabular}
\end{table}

\subsection{Assets}

The results in this draft cover five core assets. Petrobras (PETR4)
and Vale (VALE3) are the two largest companies on B3 by trading
volume. Itaú Unibanco (ITUB4) and Bradesco (BBDC4) represent the
financial sector. The Ibovespa index (\textasciicircum BVSP) serves
as the market reference. Daily Close fixtures for these five assets
are committed in the repository (\texttt{benchmark/prices/paper/})
and are the sole input of \texttt{make reproduce}. Apple (AAPL) and
the S\&P 500 ETF (SPY) are included in rule-based experiments for
cross-market comparison (Section~4.6). A broader nine-asset universe
adding mid-cap and sector-diverse tickers (EGIE3, SLCE3, RENT3,
RADL3) is defined in the benchmark for future agent evaluations; no
results on those four tickers are reported here.

\subsection{Agents}

All agents implement a uniform interface: they receive the same
\texttt{market\_state} object --- containing OHLCV prices and
macroeconomic indicators (SELIC, IPCA, USD/BRL) --- and return one
of three signals: \texttt{BUY}, \texttt{SELL}, or \texttt{HOLD}.
Rule-based agents extract only price data from \texttt{market\_state};
the macro context is carried in the interface so that macro-aware
agents (e.g.\ LLM agents) can be plugged in without changing the
simulation loop. No macro-aware agent is evaluated in this draft.

\begin{table}[h]
\centering
\caption{Agents evaluated in BrazilBench.}
\label{tab:agents}
\begin{tabular}{lp{8cm}}
\toprule
\textbf{Agent} & \textbf{Decision Logic} \\
\midrule
Buy-and-Hold  & BUY on day 1, HOLD thereafter. Passive baseline. \\
MACD          & MACD(12,26,9) crossover: BUY on upward cross,
                SELL on downward cross. \\
SMA Crossover & SMA50/SMA200 golden/death cross. Requires 200-day
                warmup. \\
Momentum      & BUY if 60d return $>0\%$; SELL if $<-5\%$; else HOLD. \\
Random        & Weighted random: P(BUY)=0.3, P(SELL)=0.1, P(HOLD)=0.6;
                seed=42 (fixed for reproducibility; represents a
                single realization of the random process --- multi-seed
                Random is left for future work; Table~\ref{tab:random-n100}
                reports the $N{=}100$-seed distribution). \\
\bottomrule
\end{tabular}
\end{table}

\subsection{Scope of the Zero-Cost Reproducible Artifact}
\label{sec:scope}

Two of the result tables in this paper are regenerated verbatim by
\texttt{make reproduce} from the committed Close fixtures, with no
API key, no network access, and no GPU: the RandomAgent null
distribution over $N{=}100$ seeds (Table~\ref{tab:random-n100},
\texttt{docs/paper\_random\_n100.tex}) and the equal-weight
portfolio baselines (Table~\ref{tab:ew-baselines},
\texttt{docs/paper\_ew\_portfolio\_baselines.tex}).
\texttt{tests/test\_reproduce.py} pins both outputs byte-for-byte in
CI, and \texttt{make docker-bench} runs the same target inside a
container built from the locked environment.

The remaining classical tables (Tables~\ref{tab:ibov},
\ref{tab:return_bvsp}, \ref{tab:sharpe_petr4}, \ref{tab:cross_market})
and the figures were computed from an earlier download of the same
yfinance series and are typed into the paper; \texttt{make reproduce}
does not regenerate them. They are retained as classical results but
should be read with that caveat.

No LLM agent is part of the artifact evaluated here. The agent
interface accepts macro context so that LLM-based agents can be
evaluated on the same regimes, but any such evaluation requires paid
cloud inference or local GPU inference, lies outside the zero-cost
reproduce path, and is not claimed in this draft
(Section~\ref{sec:limitations}).

\subsection{Simulation Engine}

The simulation runs day-by-day, with each agent making one decision per
trading day. A 250-day warmup window precedes each regime, providing
slow indicators such as SMA200 with sufficient historical data before
decisions begin. Transaction costs are set at 25 basis points per fill,
reflecting realistic B3 execution costs.

\subsection{Evaluation Metrics}

We report six metrics: Total Return, Sharpe Ratio, Sortino Ratio,
Maximum Drawdown, Calmar Ratio, and Win Rate. All risk-adjusted metrics
use a risk-free rate of 11.75\% annually, corresponding to the average
SELIC rate during 2019--2022.

% -------------------------------------------------------
\section{Results}
% -------------------------------------------------------

\subsection{Market Benchmark}

The COVID-19 crisis regime produced the most severe drawdown of the four
periods, with the Ibovespa falling 45.5\% peak-to-trough. The
\texttt{high\_rates\_2022} regime presents a counterintuitive result:
the index finished the year with a positive nominal return of $+5.6\%$,
yet delivered a negative Sharpe ratio of $-0.19$. This occurs because
the SELIC risk-free rate of 11.75\% exceeded the realized return,
meaning a passive investor would have been better off holding government
bonds.

\begin{table}[h]
\centering
\caption{\^{}BVSP (Ibovespa) Buy-and-Hold --- market benchmark across regimes.}
\label{tab:ibov}
\begin{tabular}{lrrrr}
\toprule
\textbf{Regime} & \textbf{Total Return} & \textbf{Sharpe}
  & \textbf{Sortino} & \textbf{Max DD} \\
\midrule
bull\_2019     & +27.1\% &  0.80 &  1.12 & $-10.0\%$ \\
crisis\_2020   & $-23.9\%$ & $-1.00$ & $-1.23$ & $-45.5\%$ \\
recovery\_2021 & +6.7\%  &  0.18 &  0.23 & $-12.0\%$ \\
high\_rates\_2022     & +5.6\%  & $-0.19$ & $-0.29$ & $-20.9\%$ \\
\bottomrule
\end{tabular}
\end{table}

\subsection{Strategy Performance by Regime}

Table~\ref{tab:return_bvsp} reports total return per agent on
\textasciicircum BVSP across the four regimes.
Table~\ref{tab:sharpe_petr4} reports Sharpe ratio per agent on PETR4.
No single strategy dominates across all four regimes, motivating the
multi-regime design of BrazilBench.

The Random agent uses a fixed seed (seed=42) throughout the benchmark
for reproducibility. To characterize the full null distribution,
Table~\ref{tab:random-n100} reports the distribution of RandomAgent
returns over $N{=}100$ seeds per cell. The seed=42 draw sits at
median percentile 37.0 (mean 36.3) across 20 cells --- below-typical
but not systematically extreme. It falls in a tail ($<$5th or $>$95th
pctile) in 4 of 20 cells, most severely in crisis\_2020 PETR4 (0.5th
pctile, $-57.6\%$). This quantifies the noise floor of an uninformed
policy and confirms that single-seed Random is a noisy but broadly
representative point estimate. Equal-weight portfolio comparisons use
the $N{=}100$ mean rather than seed=42 (Table~\ref{tab:ew-baselines},
Appendix~\ref{app:portfolio}).

\begin{table}[t]
\centering
\small
\begin{tabular}{llrrrrr}
\toprule
Regime & Ticker & Mean$\pm$SD (\%) & P5 (\%) & P95 (\%) & seed42 (\%) & seed42 pctile \\
\midrule
\multirow{5}{*}{bull\_2019} & PETR4 & $15.7\pm16.0$ & $-10.5$ & $43.6$ & $0.4$ & $14.5$ \\
 & VALE3 & $6.2\pm19.4$ & $-24.6$ & $40.8$ & $22.0$ & $80.5$ \\
 & ITUB4 & $4.2\pm13.1$ & $-14.2$ & $28.2$ & $0.9$ & $40.5$ \\
 & BBDC4 & $8.6\pm15.0$ & $-13.0$ & $32.7$ & $8.6$ & $53.5$ \\
 & IBOV  & $13.9\pm9.9$  & $0.2$   & $29.5$ & $10.0$ & $39.5$ \\
\midrule
\multirow{5}{*}{crisis\_2020} & PETR4 & $-17.2\pm24.9$ & $-48.9$ & $25.0$ & $-57.6$ & $0.5$ \\
 & VALE3 & $1.0\pm19.4$   & $-27.8$ & $31.5$ & $-22.3$ & $13.5$ \\
 & ITUB4 & $-13.0\pm15.3$ & $-34.2$ & $12.7$ & $-20.9$ & $34.5$ \\
 & BBDC4 & $-19.1\pm16.3$ & $-38.8$ & $9.1$  & $-33.7$ & $19.5$ \\
 & IBOV  & $-13.1\pm15.4$ & $-34.0$ & $15.4$ & $-33.6$ & $6.5$ \\
\midrule
\multirow{5}{*}{recovery\_2021} & PETR4 & $5.9\pm18.8$  & $-21.7$ & $38.7$ & $-21.3$ & $7.5$ \\
 & VALE3 & $15.3\pm12.8$ & $-3.5$  & $35.0$ & $27.8$  & $83.5$ \\
 & ITUB4 & $0.2\pm12.3$  & $-17.3$ & $22.8$ & $-28.2$ & $0.5$ \\
 & BBDC4 & $5.3\pm12.2$  & $-11.2$ & $26.8$ & $-20.4$ & $0.5$ \\
 & IBOV  & $4.4\pm6.8$   & $-5.8$  & $15.4$ & $-5.9$  & $4.5$ \\
\midrule
\multirow{5}{*}{hi\_rates\_2022} & PETR4 & $15.2\pm25.4$ & $-21.9$ & $62.6$ & $9.6$  & $45.5$ \\
 & VALE3 & $11.8\pm21.7$ & $-19.7$ & $45.0$ & $5.5$  & $43.5$ \\
 & ITUB4 & $10.2\pm16.7$ & $-13.2$ & $40.4$ & $33.6$ & $89.5$ \\
 & BBDC4 & $-7.4\pm17.3$ & $-30.1$ & $20.1$ & $7.5$  & $83.5$ \\
 & IBOV  & $1.6\pm11.3$  & $-13.2$ & $21.2$ & $4.3$  & $64.5$ \\
\bottomrule
\end{tabular}
\caption{\textbf{RandomAgent null distribution, $N{=}100$ seeds.}
  Distribution of total return of an uninformed RandomAgent (BUY/SELL/HOLD
  drawn i.i.d.\ uniformly each day) over seeds 0--99. The seed=42 draw
  sits at median percentile 37.0 (mean 36.3) across 20 cells; it falls
  in a tail ($<$5th or $>$95th pctile) in 4 of 20 cells.
  Wide spreads in crisis\_2020 reflect path dependence of random
  entry/exit timing under high volatility.
  IBOV $=$ \texttt{\^{}BVSP}. Auto-generated by
  \texttt{scripts/run\_random\_n100.py}.}
\label{tab:random-n100}
\end{table}

\begin{table}[h]
\centering
\caption{Total return by agent on \^{}BVSP across four regimes.
  \textbf{Random = seed=42 (single realization)}; this is one draw
  from the random baseline distribution, not a mean. All other agents
  are deterministic given price data.}
\label{tab:return_bvsp}
\begin{tabular}{lrrrr}
\toprule
\textbf{Agent} & \textbf{\footnotesize bull\_2019}
  & \textbf{\footnotesize crisis\_2020}
  & \textbf{\footnotesize recovery\_2021}
  & \textbf{\footnotesize hi\_rates\_2022} \\
\midrule
Buy-and-Hold  & +27.1\% & $-23.9\%$ & +6.7\%  & +5.6\%  \\
MACD          & +11.8\% & \textbf{+9.8\%} & +1.8\%  & +6.9\%  \\
SMA Crossover & 0.0\%$^\dag$ & 0.0\%$^\dag$ & 0.0\%$^\dag$ & $-14.6\%$ \\
Momentum      & +20.0\% & $-11.3\%$ & +4.1\%  & $-5.5\%$  \\
Random        & +13.2\% & $-12.4\%$ & $-7.5\%$ & $-5.5\%$  \\
\bottomrule
\multicolumn{5}{l}{\footnotesize $^\dag$ SMA Crossover: $N{=}0$ trades.
  Mean/Median not comparable to agents with $N{=}20$ cells.}
\end{tabular}
\end{table}

\begin{table}[h]
\centering
\caption{Sharpe ratio by agent on PETR4 across four regimes.
  $r_f = 11.75\%$ (SELIC).}
\label{tab:sharpe_petr4}
\begin{tabular}{lrrrr}
\toprule
\textbf{Agent} & \textbf{\footnotesize bull\_2019}
  & \textbf{\footnotesize crisis\_2020}
  & \textbf{\footnotesize recovery\_2021}
  & \textbf{\footnotesize hi\_rates\_2022} \\
\midrule
Buy-and-Hold  &  0.53 & $-0.50$ &  0.09 & $-0.39$ \\
MACD          & $-0.50$ & \textbf{1.53} & $-0.73$ & $-0.61$ \\
SMA Crossover & $-0.93$ & 0.0$^{\dag}$ &  1.24 & $-1.57$ \\
Momentum      & $-0.74$ & $-2.74$ & $-1.24$ & $-0.36$ \\
Random        & $-0.17$ & $-0.11$ & $-0.84$ & $-1.45$ \\
\bottomrule
\multicolumn{5}{l}{\footnotesize $^{\dag}$ SMA Crossover: no trades
  executed in 80-day regime.}
\end{tabular}
\end{table}

\subsection{Case Study: MACD Avoids the COVID Crash}

\textbf{Note: this section presents a single-event case study, not a
general finding.} The MACD result on PETR4/crisis\_2020 stems from
3 fills --- one well-timed exit/re-entry sequence. It should not be
interpreted as evidence of general MACD superiority.

For balance: in bull\_2019, MACD captured only $+11.8\%$ vs.\ B\&H's
$+27.1\%$ on \^{}BVSP (Table~\ref{tab:return_bvsp}) and $+19.4\%$
vs.\ B\&H's $+25.0\%$ on PETR4 --- losing to the passive baseline
by 5--15pp in a sustained uptrend. Table~\ref{tab:return_bvsp} shows this
directly: crisis\_2020 is the \emph{only} regime where MACD
outperforms B\&H. Figure~1 was chosen because it is the most
visually informative case; the full picture is less flattering to
MACD.

MACD exited PETR4 at day~8 on a downward crossover signal and
re-entered at day~37, staying in cash through the March 2020 crash.
The result was $+26.2\%$ vs.\ Buy-and-Hold's $-28.0\%$, a 54.2pp
spread. \textbf{This result requires careful interpretation: the
entire advantage comes from three fills (one exit and one re-entry
on PETR4), constituting essentially a single well-timed market-timing
decision.} It cannot be read as evidence of general MACD superiority.

The more defensible number is on \^{}BVSP: $+9.8\%$ vs.\ $-23.9\%$,
a 33.7pp spread, also from 3 fills but on the index rather than a
single stock, and with no single-stock idiosyncratic risk. The pattern
generalized across VALE3 ($+13.7\%$) and the Ibovespa, suggesting
the MACD exit signal captured the macro crash rather than any
PETR4-specific event --- but a single 80-day window is insufficient
to claim robust strategy superiority over Buy-and-Hold across regimes;
no formal significance test is reported in this draft.

\begin{figure}[h]
\centering
\includegraphics[width=0.88\linewidth]{figures/equity_curve_crisis2020_petr4.pdf}
\caption{Equity curves for MACD vs.\ Buy-and-Hold on PETR4 during
  crisis\_2020 (2020-02-03 to 2020-05-29, 80 trading days).
  MACD exits at day 8 on a downward crossover signal, stays in cash
  through the March 2020 crash, and re-enters at day 37 as momentum
  reverses. Buy-and-Hold absorbs the full drawdown. Final spread:
  $+54.2$ percentage points.}
\label{fig:equity_curve}
\end{figure}

\subsection{Finding 2 --- Momentum Whipsaw in High-Rates Regime}

The high-rates 2022 regime on B3 was characterized by high volatility
rather than a clean downtrend: the Ibovespa alternated between
commodity-driven rallies and rate-driven selloffs throughout the year.
Momentum evaluates the 60-day return and buys when positive, sells when
sufficiently negative. In a market that oscillates without sustained
direction, this causes the agent to repeatedly buy at local highs and
sell at local lows. With 8--9 round-trips at 25 basis points each,
transaction costs compound and the result is severe: Momentum lost
$-5.5\%$ on the Ibovespa and $-28.5\%$ at worst across the Brazilian
asset universe in the high-rates 2022 regime, compared to
Buy-and-Hold's $+5.6\%$ on the same index. The U.S.\ assets (AAPL,
SPY) in the cross-market validation exhibited the same whipsaw pattern,
confirming the effect is strategy-driven rather than market-specific
--- but the primary finding applies to B3 assets.

\subsection{Finding 3 --- Strong Trends Defeat Active Management}

In bull\_2019, Buy-and-Hold on the Ibovespa returned $+27.1\%$ while
MACD captured only $+11.8\%$ and Momentum $+20.0\%$. In a sustained
uptrend, any exit from the market is a missed opportunity. MACD exited
during minor corrections and paid the re-entry cost twice. The
conclusion generalizes across all assets tested: in strongly trending
markets, every active signal introduces the risk of selling too early
and buying back too late. The U.S.\ assets in the cross-market set
(AAPL $+85\%$ in bull\_2019) exhibit the same pattern at larger
magnitude, which we report in the cross-market comparison below.

\subsection{Finding 4 --- Risk-Free Rate Determines Cross-Market Comparability}

The same strategy applied over the same period produces radically
different Sharpe ratios depending on which country's risk-free rate is
applied. Table~\ref{tab:cross_market} illustrates this with Buy-and-Hold
in bull\_2019.

\begin{table}[h]
\centering
\caption{Buy-and-Hold Sharpe ratio in bull\_2019 under two risk-free
  rate conventions. Using a single rate for both markets produces
  misleading comparisons.}
\label{tab:cross_market}
\begin{tabular}{lrrrr}
\toprule
\textbf{Asset} & \textbf{Return} & \textbf{Sharpe (SELIC 11.75\%)}
  & \textbf{Sharpe (Fed Funds 2.4\%)} & \textbf{Correct $r_f$} \\
\midrule
PETR4 (BR) & +25.0\% & 0.53 & 1.68 & SELIC \\
VALE3 (BR) & +17.8\% & 0.21 & 1.09 & SELIC \\
AAPL  (US) & +85.0\% & 4.87 & 5.42 & Fed Funds \\
SPY   (US) & +28.9\% & 1.62 & 2.05 & Fed Funds \\
\bottomrule
\end{tabular}
\end{table}

Three observations follow. First, applying SELIC (11.75\%) to AAPL
would produce a Sharpe of 4.87 --- artificially low relative to a
U.S.\ investor's true opportunity cost of 2.4\%. Conversely, applying
Fed Funds (2.4\%) to PETR4 would produce 1.68 --- grossly overstating
the risk-adjusted return a Brazilian investor actually experienced.
Neither comparison is valid. Second, even under country-appropriate
rates, PETR4's Sharpe (0.53) is substantially below SPY (2.05) ---
not because of lower return, but because the higher Brazilian
risk-free rate consumes most of the nominal gain. A 25\% return in
Brazil generates only 13\% excess return; the same 25\% in the U.S.\
generates 22.6\%. Third, this asymmetry is structural, not
coincidental: it reflects the persistent interest-rate differential
between emerging and developed markets that any cross-market benchmark
must account for explicitly. All BrazilBench metrics for BR assets
use SELIC; all metrics for U.S.\ assets use the Fed Funds rate for
the corresponding year.

% -------------------------------------------------------
\section{Discussion}
% -------------------------------------------------------

\begin{figure}[h]
\centering
\includegraphics[width=0.88\linewidth]{figures/tc_sensitivity.pdf}
\caption{Strategy robustness to transaction costs (5 rule-based agents,
  5 BR tickers $\times$ 4 regimes, mean total return across all cells).
  Buy-and-Hold is nearly immune (1 fill). MACD is the only active strategy
  with positive return at zero cost, breaking even at $\sim$3\,bps.
  Momentum is the most TC-sensitive (8.7 avg fills/cell), declining from
  $-0.7\%$ at zero cost to $-9.6\%$ at 100\,bps. No strategy is robustly
  cost-efficient (positive at 100\,bps) other than Buy-and-Hold.
  Vertical dashed line: BrazilBench default of 25\,bps.}
\label{fig:tc}
\end{figure}

The SMA200 requires 200 days of price history. The crisis\_2020 regime
spans only 80 trading days, and no golden or death cross occurred within
the regime even with the 250-day warmup. This is a methodological
finding, not a failure of the indicator: SMA crossover was designed for
long sustained trends, not short crisis windows. A single-period
benchmark would obscure this entirely.

The Random agent has no logic. If an intelligent strategy cannot
consistently outperform Random across multiple regimes, it has no real
edge. Any strategy that ties with Random is operating at the same level
of randomness, regardless of how sophisticated its signal logic appears.

Transaction costs reveal a critical gap between gross and net
performance. MACD on PETR4 in bull\_2019 achieved $+0.97\%$ gross
return. After adjusting for 20 fills at 25 basis points each, the
cost-adjusted return drops to $-4.0\%$. Figure~\ref{fig:tc} shows
this effect systematically: MACD is the only active strategy with
positive mean return at zero cost, but breaks even at approximately
3\,bps. Momentum (8.7 avg fills/cell) is the most TC-sensitive,
declining from $-0.7\%$ at zero cost to $-9.6\%$ at 100\,bps.
Buy-and-Hold is nearly immune (1 fill). No active strategy delivers
positive mean return at BrazilBench's default of 25\,bps.

% -------------------------------------------------------
\section{Limitations and Future Work}
\label{sec:limitations}
% -------------------------------------------------------

\begin{enumerate}

\item \textbf{LLM agent evaluation is out of scope for this draft.}
BrazilBench's agent interface carries macro context so that LLM-based
agents --- cloud APIs or open-weight models --- can be evaluated on the
same regimes and assets. This draft reports no such results: LLM
inference is paid or GPU-bound, is not part of \texttt{make
reproduce}, and is therefore not claimed here. A future LLM
evaluation should, at minimum, (a) report $N{\geq}3$ independent
sessions per cell with variance, since \texttt{temperature=0} does
not guarantee determinism in production inference; (b) test
predictive superiority against Buy-and-Hold with the Hansen SPA test
\cite{hansen2005} and the Diebold-Mariano test \cite{dieboldmariano1995}
under multiple-comparison correction \cite{harveyliuzhu2016}; (c)
address pretraining-memory leakage, since all four regimes
(2019--2022) predate the training cutoff of current models, e.g.\ via
date-shuffle probes and post-cutoff out-of-sample regimes; and (d)
ablate the macro features, which rule-based agents do not consume.

\item \textbf{Hand-typed classical tables.} Tables~\ref{tab:ibov},
\ref{tab:return_bvsp}, \ref{tab:sharpe_petr4} and \ref{tab:cross_market}
and the figures come from an earlier price vintage and are not
regenerated by \texttt{make reproduce} (Section~\ref{sec:scope}).
Only Tables~\ref{tab:random-n100} and \ref{tab:ew-baselines} are
CI-pinned. Migrating the remaining tables onto the fixture-based
path is the immediate next step.

\item \textbf{Regime definition robustness.} The four benchmark regimes
are defined by hand based on macroeconomic events. While this makes
them interpretable, it creates exposure to ex-post selection:
a researcher who saw the data first could define regimes that favor
one strategy. A robustness check using Hamilton (1989) Hidden Markov
Model \cite{hamilton1989} regime detection on the same period would
test whether conclusions are stable across regime definitions. Visual
inspection of HMM-detected breakpoints on the Ibovespa (not formally
reported) suggests broad alignment with the hand-defined dates, but
formal documentation and sensitivity analysis should be included in
a future version.

\item \textbf{Single-asset simulation.} Each run allocates 100\% of
capital to a single asset. Appendix~\ref{app:portfolio} provides an
equal-weight portfolio estimate using existing data; full multi-asset
simulation with correlation structure is left for future work.

\item \textbf{Data source limitations.} Historical prices are sourced
from yfinance with the \texttt{.SA} suffix for Brazilian assets.
yfinance is adequate for reproducible academic benchmarks but has
known limitations in production: gaps around Brazilian-specific
holidays, inconsistent dividend adjustment methodology, and occasional
errors in split ratios for B3-listed companies. These issues can
introduce price discontinuities of 0.1--2\% on affected dates.
For a production-grade benchmark, data from B3 directly or from
commercial providers such as Economatica or Refinitiv Eikon would
eliminate these inconsistencies. We chose yfinance to maximize
reproducibility and recommend validation against a commercial source
before any live trading application.

\item \textbf{Survivorship bias.} We tested only liquid, continuously
traded assets on B3. Companies that failed or were delisted during the
benchmark period are absent. \textbf{Sensitivity bound:}
\cite{lesmond1999} estimates 1--3pp annual upward distortion for
liquid-only sets in U.S.\ emerging market data of the 1990s; the
applicability to Brazilian equities of 2019--2022 is uncertain.
Such a correction would inflate absolute Sharpe ratios but applies
equally to all agents on the same asset. Future work should include
distressed assets (OIBR3, GOLL4, MGLU3) to bracket the effect
empirically rather than relying on a cross-decade, cross-country
estimate.

\item \textbf{Random baseline: single realization.} The Random agent
uses a fixed seed (seed=42) for reproducibility. This produces one
specific realization of the random signal process, not an expected
value over the random baseline distribution. Table~\ref{tab:random-n100}
reports the full $N{=}100$-seed distribution: seed=42 sits at median
percentile 37.0 and falls in a distribution tail in 4 of 20 cells.
Single-seed Random is therefore a noisy but broadly representative
estimate. Equal-weight portfolio comparisons (Table~\ref{tab:ew-baselines})
use the $N{=}100$ mean instead of seed=42.

\item \textbf{Asset-specific slippage.} The benchmark uses 25\,bps
for all tickers. A sensitivity analysis with liquidity premia
(EGIE3/RADL3 $+5$\,bps; SLCE3/RENT3 $+10$\,bps) produces mean
return deltas of $-0.25$\,pp and $-0.51$\,pp for mid-cap buckets
--- small but non-negligible for high-turnover strategies. The
realistic model is available via \texttt{slippage\_bps\_by\_ticker}.

\item \textbf{No formal significance testing.} The classical regime
findings are descriptive. With 80--250 trading days per regime and a
handful of fills per cell, the per-regime comparisons are
under-powered for formal tests; the case study in Section~4.3 rests
on three fills. The SPA/DM machinery listed above for future LLM
evaluation applies equally to the rule-based agents.

\end{enumerate}

% -------------------------------------------------------
\section{Conclusion}
% -------------------------------------------------------

BrazilBench is a reproducible benchmark for trading agents on
Brazilian equity markets: four macroeconomic regimes, five core B3
assets, a uniform agent interface with macro context, SELIC as the
risk-free rate, and a zero-cost reproduce path that regenerates the
RandomAgent null distribution and the equal-weight portfolio
baselines byte-for-byte in CI.

Three grouped findings, all on classical agents: \textit{(i) Regime
sensitivity of rule-based strategies:} MACD avoided the COVID crash
via a single well-timed exit/re-entry (3 fills; PETR4 $+26.2\%$ vs.\
B\&H $-28.0\%$; the more generalizable number is \^{}BVSP $+9.8\%$
vs.\ $-23.9\%$); Momentum fails under transaction costs in volatile
regimes; no active strategy consistently beats Buy-and-Hold in
sustained bull trends, and none delivers positive mean return at
25\,bps. \textit{(ii) The noise floor of an uninformed policy is
wide:} the $N{=}100$-seed Random distribution has per-cell standard
deviations of roughly 10--25\,pp, and the canonical seed=42 draw
lands in a distribution tail in 4 of 20 cells, so single-seed
baselines should not be read as expected values. \textit{(iii) The
risk-free rate convention is first-order:} applying SELIC versus Fed
Funds changes Sharpe ratios by a factor of 2--3 for the same return
series.

No LLM agent is evaluated in this draft. The benchmark is built so
that LLM agents can be run on the same regimes, but that evaluation
requires paid or GPU inference, is outside \texttt{make reproduce},
and is left as future work with the methodological requirements
listed in Section~\ref{sec:limitations}. These findings are bounded
by 20 evaluation cells (5 assets $\times$ 4 regimes) and by the
descriptive nature of the comparisons; absence of a formal test
should be read as \emph{no claim of significance}, not as evidence
of no effect.

The use of SELIC as the risk-free rate is not a minor detail. Using a
country-appropriate risk-free rate is the difference between a meaningful
and a misleading Sharpe ratio in any emerging-market benchmark.

% -------------------------------------------------------
\appendix
\section{Equal-Weight Portfolio Approximation}
\label{app:portfolio}
% -------------------------------------------------------

Table~\ref{tab:ew-baselines} reports equal-weight (1/5) portfolio
returns across the five BR tickers for each rule-based agent and
regime. Random uses the $N{=}100$-seed mean from
Table~\ref{tab:random-n100} rather than seed=42, giving the expected
uninformed return. This table is regenerated verbatim by
\texttt{make reproduce} (Section~\ref{sec:scope}).

\begin{table}[t]
\centering
\small
\begin{tabular}{lrrrr}
\toprule
\textbf{Agent (EW, \%)} & \textbf{bull\_2019} & \textbf{crisis\_2020}
  & \textbf{recovery\_2021} & \textbf{hi\_rates\_2022} \\
\midrule
Buy \& Hold     & $+17.2$ & $-22.4$ & $+10.0$ & $+14.6$ \\
MACD(12,26,9)   & $+4.9$  & $\mathbf{+6.9}$ & $+8.8$ & $+14.6$ \\
SMA(50/200)     & $+11.2$ & $-26.0$ & $+8.0$  & $-11.0$ \\
Momentum(60)    & $+4.5$  & $-4.5$  & $+3.1$  & $+12.2$ \\
Random ($N{=}100$ mean) & $+9.7$ & $-12.3$ & $+6.2$ & $+6.3$ \\
\bottomrule
\end{tabular}
\caption{\textbf{Equal-weight portfolio returns, rule-based agents only.}
  Mean of per-ticker cumulative returns across PETR4, VALE3, ITUB4,
  BBDC4, \texttt{\^{}BVSP}. Long-only, all-in/all-cash, frictionless;
  indicators warm at each window's first bar. Random row $=$ $N{=}100$
  seed mean (robust estimate; not seed=42).
  MACD is the only strategy positive in crisis\_2020 ($+6.9\%$) while
  Buy\,\&\,Hold ($-22.4\%$) and SMA ($-26.0\%$) are hit hardest;
  Buy\,\&\,Hold dominates in bull\_2019 ($+17.2\%$) via pure trend.
  Auto-generated by \texttt{scripts/run\_ew\_portfolio.py}.}
\label{tab:ew-baselines}
\end{table}

% -------------------------------------------------------

\begin{thebibliography}{9}

\bibitem{hansen2005}
Hansen, P.R.
\newblock A Test for Superior Predictive Ability.
\newblock \textit{Journal of Business \& Economic Statistics},
  23(4):365--380, 2005.

\bibitem{dieboldmariano1995}
Diebold, F.X. and Mariano, R.S.
\newblock Comparing Predictive Accuracy.
\newblock \textit{Journal of Business \& Economic Statistics},
  13(3):253--263, 1995.

\bibitem{tradingagents2024}
Wang et al.
\newblock TradingAgents: Multi-Agents LLM Financial Trading Framework.
\newblock \textit{arXiv preprint arXiv:2412.20138}, 2024.

\bibitem{finagent2024}
Zhang, X., et al.
\newblock FinAgent: A Multimodal Foundation Agent for Financial Trading.
\newblock \textit{arXiv preprint arXiv:2402.18485}, 2024.

\bibitem{chronos2024}
Ansari, A.F., et al.
\newblock Chronos: Learning the Language of Time Series.
\newblock \textit{Transactions on Machine Learning Research}, 2024.
  arXiv:2403.07815.

\bibitem{backtrader2015}
Molina, D.
\newblock Backtrader: Python Backtesting Library.
\newblock \url{https://github.com/mementum/backtrader}, 2015.

\bibitem{quantconnect2013}
QuantConnect.
\newblock Lean Algorithmic Trading Engine.
\newblock \url{https://github.com/QuantConnect/Lean}, 2013.

\bibitem{hamilton1989}
Hamilton, J.D.
\newblock A New Approach to the Economic Analysis of Nonstationary Time Series
  and the Business Cycle.
\newblock \textit{Econometrica}, 57(2):357--384, 1989.

\bibitem{angbekaert2002}
Ang, A. and Bekaert, G.
\newblock International Asset Allocation with Regime Shifts.
\newblock \textit{Review of Financial Studies}, 15(4):1137--1187, 2002.

\bibitem{lesmond1999}
Lesmond, D.A., Ogden, J.P., and Trzcinka, C.A.
\newblock A New Estimate of Transaction Costs.
\newblock \textit{Review of Financial Studies}, 12(5):1113--1141, 1999.

\bibitem{adlerdumas1983}
Adler, M. and Dumas, B.
\newblock International Portfolio Choice and Corporation Finance: A Synthesis.
\newblock \textit{Journal of Finance}, 38(3):925--984, 1983.

\bibitem{kronos2025}
Shi, Y., Fu, Z., Chen, S., Zhao, B., Xu, W., Zhang, C., and Li, J.
\newblock Kronos: A Foundation Model for the Language of Financial Markets.
\newblock \textit{arXiv preprint arXiv:2508.02739}, 2025.
  (Accepted at AAAI-26; proceedings forthcoming.)
  \url{https://ojs.aaai.org/index.php/AAAI/article/view/39730}.

\bibitem{fingpt2023}
Yang et al.
\newblock FinGPT: Open-Source Financial Large Language Models.
\newblock \textit{AI4Finance Foundation}, 2023.

\bibitem{harveyliuzhu2016}
Harvey, C.R., Liu, Y., and Zhu, H.
\newblock \ldots and the Cross-Section of Expected Returns.
\newblock \textit{Review of Financial Studies}, 29(1):5--68, 2016.

\bibitem{bacen}
Banco Central do Brasil.
\newblock Sistema Gerenciador de Series Temporais (SGS).
\newblock \url{https://api.bcb.gov.br}, 2024.

\bibitem{lopezdeprado2018}
López de Prado, M.
\newblock \textit{Advances in Financial Machine Learning}.
\newblock Wiley, 2018.

\end{thebibliography}

\end{document}
